\documentclass[aps,prb,onecolumn,nofootinbib,superscriptaddress]{revtex4-2}
\usepackage{amsmath,amssymb,amsthm,mathtools,bm}
\usepackage{booktabs}
\usepackage[colorlinks=true,linkcolor=blue,citecolor=blue,urlcolor=blue]{hyperref}
\usepackage{microtype}

\newcommand{\Id}{\mathbf{1}}
\newcommand{\dd}{\mathrm{d}}
\newcommand{\Tr}{\operatorname{Tr}}

\begin{document}

\title{Rank-One Physical Covariance Updates in the Markov Circuit}
\author{}
\date{\today}

\begin{abstract}
This note records the three algebraically equivalent forms of the rank-one
Gaussian measurement update used by the CPU and GPU Markov-circuit engines:
the dense Schur-complement solve, the rank-one resolvent action, and the
expanded Sherman--Morrison formula.  It also records the small-system survival
checks used when switching the production physical covariance update to the
rank-one resolvent implementation.
\end{abstract}

\maketitle
\setcounter{tocdepth}{2}
\tableofcontents

\section{Setup}

Let \(G\) denote the Hermitian complex-fermion covariance on the active top
layer, and let
\begin{equation}
    P = \chi \chi^\dagger,\qquad Q=\Id-P,\qquad \chi^\dagger\chi=1.
\end{equation}
We write a selected measurement branch using the sign
\begin{equation}
    s =
    \begin{cases}
        +1, & \text{occupied / particle branch},\\
        -1, & \text{unoccupied / hole branch}.
    \end{cases}
\end{equation}
The target covariance on the measured mode is therefore \(sP\).  The singular
surface of the conditional update is
\begin{equation}
    1+s\,\chi^\dagger G\chi = 0,
\end{equation}
which is the zero-probability conditioning surface for the selected branch.

\section{Dense Schur-Complement Solver}

The original implementation used the dense Gaussian contraction
\begin{equation}
    G' = H_{22} - H_{21}(G H_{11}-\Id)^{-1}G H_{21}^\dagger ,
\end{equation}
with
\begin{equation}
    (H_{11},H_{21},H_{22}) =
    \begin{cases}
        (-P,Q,P), & s=+1,\\
        (P,Q,-P), & s=-1.
    \end{cases}
\end{equation}
Equivalently,
\begin{equation}
    \boxed{
    G' = sP + Q(\Id+sGP)^{-1}GQ .
    }
    \label{eq:dense}
\end{equation}
The dense implementation evaluates Eq.~\eqref{eq:dense} by solving
\begin{equation}
    (\Id+sGP)X = GQ.
\end{equation}
In the production code this was previously routed through
\(\texttt{\_solve\_regularized}\) or \(\texttt{\_solve\_regularized\_batched}\).

\section{Rank-One Resolvent Solver}

Since \(P=\chi\chi^\dagger\),
\begin{equation}
    GP = (G\chi)\chi^\dagger .
\end{equation}
Thus the matrix \(\Id+sGP\) differs from the identity by rank one.  Defining
\begin{equation}
    v=G\chi,\qquad d_s = 1+s\chi^\dagger v,
\end{equation}
the inverse action on any matrix right-hand side \(B\) is
\begin{equation}
    \boxed{
    (\Id+sGP)^{-1}B
    =
    B-\frac{s}{d_s}v(\chi^\dagger B).
    }
    \label{eq:rankoneaction}
\end{equation}
The production CPU and GPU Markov-circuit engines now use
Eq.~\eqref{eq:rankoneaction} for the physical covariance update.  If
\(|d_s|\) is numerically unsafe, the implementation falls back to the same
small ridged dense solve convention used elsewhere in the Markov engine.

\section{Expanded Sherman--Morrison Form}

Substituting Eq.~\eqref{eq:rankoneaction} into Eq.~\eqref{eq:dense} gives the
fully expanded formula
\begin{equation}
    \boxed{
    G'
    =
    sP + QGQ
    -
    \frac{s}{1+s\chi^\dagger G\chi}
    (QG\chi)(\chi^\dagger GQ).
    }
    \label{eq:sherman}
\end{equation}
This is algebraically identical to the rank-one resolver, but it forms the
outer-product correction directly.  It is useful as a transparent reference
implementation, while the production code keeps the resolver-action form so the
support-support and support-complement right-hand sides can be handled uniformly.

\section{Support-Reduced Form}

For finite Wannier support, decompose the top-layer indices into the measured
support \(S\) and complement \(R\).  With
\begin{equation}
    G =
    \begin{pmatrix}
        G_{SS} & G_{SR}\\
        G_{RS} & G_{RR}
    \end{pmatrix},
\end{equation}
and \(P=\chi_S\chi_S^\dagger\) supported only on \(S\), the rank-one action is
applied only to the support block,
\begin{equation}
    Y=(\Id+sG_{SS}P)^{-1}G_{SR},\qquad
    Z=(\Id+sG_{SS}P)^{-1}G_{SS}Q.
\end{equation}
The updated blocks are
\begin{align}
    G'_{SS} &= sP + QZ,\\
    G'_{SR} &= QY,\\
    G'_{RR} &= G_{RR}-s\,G_{RS}PY.
\end{align}
This is the form used by the CPU and GPU local-support kernels.

\section{Small-System Survival Checks}

The diagnostic scripts use \(N_x=4\), \(N_y=8\), \(50\) cycles, one sample,
finite Wannier support \(n_{\rm shell}=1\), slab-only domain-wall dynamics, and
a fixed deterministic initial covariance.  Blowup is defined as a nonfinite
entry or \(\max_{ij}|G_{ij}|>10^{12}\).

\subsection{CPU engine}

\begin{center}
\begin{tabular}{llcccccc}
\toprule
protocol & solver & survived & first bad & final max & run max & final invol. & final spectrum\\
\midrule
postselect & rank1 & yes & -- & \(1.000\) & \(1.105\) & \(1.818\times10^{-6}\) & \([-1.000,1.000]\)\\
postselect & dense & yes & -- & \(1.000\) & \(4.561\) & \(4.670\times10^{-6}\) & \([-1.000,1.000]\)\\
postselect & Sherman & yes & -- & \(1.000\) & \(1.879\) & \(1.299\times10^{-6}\) & \([-1.000,1.000]\)\\
perfect correction & rank1 & yes & -- & \(1.000\) & \(2.132\) & \(5.701\) & \([-1.378,2.511]\)\\
perfect correction & dense & yes & -- & \(1.000\) & \(2.132\) & \(5.701\) & \([-1.378,2.511]\)\\
perfect correction & Sherman & yes & -- & \(1.000\) & \(2.132\) & \(5.701\) & \([-1.378,2.511]\)\\
\bottomrule
\end{tabular}
\end{center}

Final-state differences relative to production rank-one were
\begin{center}
\begin{tabular}{llcc}
\toprule
protocol & reference & Frobenius difference & max-entry difference\\
\midrule
postselect & dense & \(1.571\times10^{-6}\) & \(9.370\times10^{-8}\)\\
postselect & Sherman & \(2.614\times10^{-7}\) & \(2.303\times10^{-8}\)\\
perfect correction & dense & \(3.854\times10^{-14}\) & \(4.772\times10^{-15}\)\\
perfect correction & Sherman & \(4.618\times10^{-14}\) & \(8.132\times10^{-15}\)\\
\bottomrule
\end{tabular}
\end{center}

\subsection{GPU class on CPU device}

CUDA was not visible in the local environment, so the GPU-class diagnostic was
run on its CPU device path with the same Torch kernels and canonical
\texttt{classA\_U1FGTN\_gpu.run\_markov\_circuit(...)} entry point.

\begin{center}
\begin{tabular}{llcccccc}
\toprule
protocol & solver & survived & first bad & final max & run max & final invol. & final spectrum\\
\midrule
postselect & rank1 & yes & -- & \(1.000\) & \(1.000\) & \(2.344\times10^{-6}\) & \([-1.000,1.000]\)\\
postselect & dense & yes & -- & \(1.000\) & \(1.000\) & \(2.624\times10^{-6}\) & \([-1.000,1.000]\)\\
postselect & Sherman & yes & -- & \(1.000\) & \(1.000\) & \(3.449\times10^{-6}\) & \([-1.000,1.000]\)\\
perfect correction & rank1 & yes & -- & \(1.000\) & \(1.000\) & \(2.154\times10^{-13}\) & \([-1.000,1.000]\)\\
perfect correction & dense & yes & -- & \(1.000\) & \(1.000\) & \(1.920\times10^{-13}\) & \([-1.000,1.000]\)\\
perfect correction & Sherman & yes & -- & \(1.000\) & \(1.000\) & \(2.327\times10^{-13}\) & \([-1.000,1.000]\)\\
\bottomrule
\end{tabular}
\end{center}

Final-state differences relative to production rank-one were
\begin{center}
\begin{tabular}{llcc}
\toprule
protocol & reference & Frobenius difference & max-entry difference\\
\midrule
postselect & dense & \(4.717\times10^{-7}\) & \(3.326\times10^{-8}\)\\
postselect & Sherman & \(6.510\times10^{-7}\) & \(6.167\times10^{-8}\)\\
perfect correction & dense & \(5.519\times10^{-14}\) & \(7.690\times10^{-15}\)\\
perfect correction & Sherman & \(1.822\times10^{-14}\) & \(2.456\times10^{-15}\)\\
\bottomrule
\end{tabular}
\end{center}

\end{document}
